\documentclass[11pt]{article}

\usepackage[a4paper,margin=1in]{geometry}
\usepackage{amsmath,amssymb,amsthm,mathtools}
\usepackage{algorithm}
\usepackage{algpseudocode}
\usepackage{hyperref}
\usepackage{enumitem}
\usepackage{booktabs}
\usepackage{graphicx}
\usepackage{float}

\hypersetup{colorlinks=true,linkcolor=blue,citecolor=blue,urlcolor=blue}

\title{A Monotone Single-Pointer Anchor-Advancement Algorithm\\
for Cycle Detection: Exact Closed Form, Provable Dominance,\\
and Optimality in the Successor-Evaluation Model}
\author{zx l}
\date{June 2026}

\newtheorem{theorem}{Theorem}[section]
\newtheorem{lemma}[theorem]{Lemma}
\newtheorem{proposition}[theorem]{Proposition}
\newtheorem{corollary}[theorem]{Corollary}
\newtheorem{definition}[theorem]{Definition}
\newtheorem{remark}[theorem]{Remark}

\newcommand{\nextop}{\texttt{next}}
\newcommand{\nullnode}{\texttt{null}}
\newcommand{\oldc}{\texttt{old}}
\newcommand{\newc}{\texttt{newc}}

\begin{document}

\maketitle

\begin{abstract}
Given a start node and a successor oracle $f$, detecting the eventual cycle of
the sequence $x,f(x),f^2(x),\dots$ and recovering its tail length $\mu$ and
period $\lambda$ is the textbook problem solved in $O(1)$ memory by Floyd's
tortoise-and-hare and Brent's power-of-two algorithms. We present a
\emph{monotone single-pointer anchor-advancement} algorithm for the same
problem. It keeps a current anchor, a candidate next anchor, and a safe
localization anchor; it probes in two phases and applies a failure-compression
(``nextval'') rule that never rescans an interval whose return-to-anchor checks
have already failed. The pointer never rewalks a node. Under the
successor-evaluation cost model we give a single, coherent account of the method:
\textbf{(i)} a correctness proof (linear time, constant space, safe
localization); \textbf{(ii)} an \emph{exact closed form} for the
evaluation count $A(\mu,\lambda)=2(\mu+\lambda)+R$ with overhead
$R\in\{1,2^{r_{\min}}-1\}$, verified bit-for-bit over all $4{\times}10^{6}$ cells
with $\mu,\lambda\le 2000$; \textbf{(iii)} a \emph{proved} per-instance dominance
theorem with a closed-form tie set --- $A\le B$ with equality \emph{iff}
$(\mu,\lambda)\in\{(0,1),(1,1),(1,2)\}$, and $A\le F$ with equality only at
$(1,1)$; \textbf{(iv)} a matching $\mu+\lambda$ lower bound proving all three
methods constant-factor optimal, with $A$ within a factor $3$ of the floor on
every input; \textbf{(v)} an \emph{average-case} constant
$c_A=\mathbb E[A]/\mathbb E[\mu+\lambda]\approx 2.18033$, shown by Mellin analysis
to be a \emph{log-periodic} function of $n$, against $\approx 3.08$ (Brent) and
$\approx 3.5$ (Floyd); \textbf{(vi)} a multi-anchor extension tracing a
detection-latency / total-work Pareto frontier; and \textbf{(vii)} wall-clock and
Pollard-$\rho$ benchmarks locating the crossover at which fewer successor
evaluations become a real speedup. We are deliberately precise about scope: the
asymptotic class $\Theta(\mu+\lambda)$ is unchanged, the method performs more
pointer comparisons than the classics, and the win is realized when a single
successor evaluation is expensive --- the Pollard-$\rho$ regime that originally
motivated Brent.
\end{abstract}

\section{Introduction}
\label{sec:intro}

Many computations reduce to iterating a deterministic map under three
constraints: the map is a \emph{black box} (only $y=f(x)$ is observable, not its
global structure); we may take only \emph{local steps} (apply $f$); and we have
\emph{bounded memory} (we cannot store the trajectory). Under these constraints a
trajectory $x_0,x_1=f(x_0),\dots$ over a finite state space is eventually
periodic: a tail of length $\mu\ge 0$ leads into a cycle of length $\lambda\ge 1$
(the ``$\rho$'' shape). Recovering $(\mu,\lambda)$ is the core primitive behind
Pollard's $\rho$ factorization and discrete-log methods, period testing of
pseudo-random generators, non-termination and lasso detection in model checking,
limit-cycle detection in iterated dynamics, and routing-loop or
stream-deduplication checks. The same problem is the classical ``find the entry
of a cycle in a singly linked list,'' with the \nextop{} pointer as the
successor oracle. In all of these settings the successor is the unit of cost, and
in several of them (cryptographic maps, heavy numerical iterations) a single
successor evaluation dominates the running time.

Floyd's tortoise-and-hare (popularized by Knuth) and Brent's power-of-two
blocking algorithm both solve the problem in $O(1)$ memory but spend,
respectively, about $3.5$ and $3.08$ successor evaluations per unit of
$\mu+\lambda$. The information-theoretic floor in the local-oracle model is
$\mu+\lambda$ evaluations: every reachable distinct node must be touched at least
once (Theorem~\ref{thm:lowerbound}). This paper closes part of that gap with a
method whose average constant is $\approx 2.18$, which provably dominates both
classics on \emph{every} input, and which --- unlike its predecessors --- admits
an exact, verifiable closed-form analysis.

The algorithm is inspired by the following observation. In Brent-style anchor
methods, after a failed round the first phase of the next round may repeat
exactly the interval already scanned as the second phase of the failed round. If
that repeated interval was already checked for returning to the relevant anchor,
rescanning it cannot produce a new outcome; we therefore \emph{compress} the
failed interval. The method maintains three distinguished pointers:
\begin{itemize}[leftmargin=2em]
    \item $c$, the current anchor;
    \item $\newc$, the candidate next anchor produced by the first phase of a round;
    \item $\oldc$, a localization anchor proved to lie before the cycle entry
          (or to be the entry itself in the boundary case $\mu=0$).
\end{itemize}
The localization anchor is the second source of savings: because $\oldc$ is in
general strictly closer to the entry than the head, entry localization is
strictly cheaper than Brent's.

Let $A(\mu,\lambda)$, $B(\mu,\lambda)$, and $F(\mu,\lambda)$ denote the number of
successor evaluations performed by the proposed, Brent's, and Floyd's algorithms
on a cyclic input with tail $\mu$ and period $\lambda$. We prove
$A\le B$ and $A\le F$ on every cyclic input, give the exact value of $A$, and
characterize precisely when equality holds.

\paragraph{Scope of the claim.}
This is not a claim of universal wall-clock superiority. The method performs
additional pointer comparisons and maintains additional state; its advantage is
stated under the \emph{successor-evaluation model}, in which each operation
$p\leftarrow f(p)$ has unit cost and comparisons, assignments, and integer
arithmetic are free. This is the natural model when successor evaluation
dominates, as in functional iteration, Pollard-$\rho$ factorization, or
pointer-machine settings. Sections~\ref{sec:wallclock}--\ref{sec:rho} quantify,
on real hardware, exactly how expensive a successor step must be for the
model-level advantage to become a wall-clock advantage.

\paragraph{Contributions and roadmap.}
\begin{enumerate}[leftmargin=2em,itemsep=2pt]
    \item A constant-space, monotone single-pointer anchor-advancement algorithm
          (Section~\ref{sec:algo}), with a probe-schedule invariant proving the
          skip rule misses no successful return-to-anchor event
          (Section~\ref{sec:schedule}) and a full correctness proof
          (Section~\ref{sec:correctness}).
    \item An \emph{exact closed form} for the evaluation count
          (Theorem~\ref{thm:closed-form}), verified over $4{\times}10^{6}$ cells
          (Section~\ref{sec:closedform}).
    \item A \emph{proved} tight dominance result over Brent and Floyd with a
          closed-form tie set (Section~\ref{sec:dominance}), refining the earlier
          ``strict on infinitely many inputs'' statement to a finite, explicit
          equality set.
    \item A matching $\mu+\lambda$ lower bound and constant-factor optimality
          (Section~\ref{sec:lowerbound}).
    \item An average-case analysis: $c_A\approx 2.18033$ as a log-periodic
          function of $n$ via Mellin residues, validated against Monte-Carlo
          (Section~\ref{sec:average}).
    \item A multi-anchor space--time tradeoff tracing a detection / total-work
          Pareto frontier (Section~\ref{sec:tradeoff}).
    \item Reproducible operation-count, wall-clock, and Pollard-$\rho$ evaluations
          locating the cheap/expensive-successor crossover
          (Sections~\ref{sec:empirical}--\ref{sec:rho}).
\end{enumerate}

\section{Model and Notation}
\label{sec:model}

Let $f:S\to S$ on a finite set $S$ with $|S|=n$, and let the sequence generated
from the start $x_0$ by following $f$ (equivalently, \nextop{} pointers) be
$v_0,v_1,v_2,\dots$. There exist unique integers $\mu\ge 0$ (tail length) and
$\lambda\ge 1$ (period) with
\[
v_{\mu+k+\lambda}=v_{\mu+k}\qquad\text{for all }k\ge 0,
\]
and $v_i\ne v_j$ for distinct $i,j<\mu+\lambda$. We identify each absolute
position $x\in\mathbb N$ with the node
\[
N(x)=\begin{cases}x,&x<\mu,\\ \mu+((x-\mu)\bmod \lambda),&x\ge \mu,\end{cases}
\]
so that positions $x,y$ denote the same node iff $N(x)=N(y)$. Write
$k=\mu+\lambda$.

\begin{definition}[Cost model]
\label{def:cost}
The cost of a computation is the number of successor evaluations
$p\leftarrow f(p)$. Pointer comparisons, assignments, Boolean tests, and integer
arithmetic are not counted. Dominance statements are made only for cyclic inputs;
for acyclic inputs we prove only linear-time termination.
\end{definition}

\begin{remark}[Absolute positions versus pointer values]
The analysis uses a lifted execution trace in which each successor evaluation
advances the moving pointer by one absolute position. In a cyclic structure,
different absolute positions may represent the same concrete node, so a stored
pointer value does not uniquely identify a position, but the chronological trace
does. All interval statements below refer to this lifted trace.
\end{remark}

\section{The Algorithm}
\label{sec:algo}

The algorithm operates in rounds. A normal round with current anchor $c$ and step
length $s$ (a power of two) has two phases:
\begin{enumerate}[leftmargin=2em]
    \item start from $c$, advance $s$ steps, and check whether the probe returns
          to $c$;
    \item let the endpoint of the first phase be $\newc$; advance $2s$ further
          steps and check whether the probe returns to either $c$ or $\newc$.
\end{enumerate}
A return to $c$ yields the cycle length directly; a return to $\newc$ yields the
cycle length measured from $\newc$ and stores the previous $c$ as $\oldc$. If the
round fails, we set $c\leftarrow\newc$, double $s$, and \emph{skip} the next first
phase because it would exactly repeat the failed second-phase interval --- reusing
the already-reached node as the start of the next round, so no node is ever
walked twice. Once $\lambda$ is known, a standard offset-and-lockstep pass
(\textsc{LocateEntry}) recovers the entry and $\mu$.

\begin{algorithm}[p]
\footnotesize
\caption{Monotone Single-Pointer Anchor Advancement}
\label{alg:nextval}
\begin{algorithmic}[1]
\Require head / start node, successor oracle $f$
\Ensure cycle entry, or \nullnode{} if the structure is acyclic
\If{\(\textit{head}=\nullnode\)} \State \Return \nullnode \EndIf
\State \(\oldc\gets \textit{head}\); \(c\gets \textit{head}\); \(s\gets 1\); \(\textit{skip}\gets \textbf{false}\)
\While{\textbf{true}}
    \If{\(\textit{skip}\)}
        \State \(p\gets \textit{skipEnd}\); \(\textit{skip}\gets \textbf{false}\)
    \Else
        \State \(p\gets c\)
        \For{\(i=1\) to \(s\)}
            \If{\(f(p)=\nullnode\)} \State \Return \nullnode \EndIf
            \State \(p\gets f(p)\)
            \If{\(p=c\)} \State \Return \Call{LocateEntry}{$\oldc,i$} \EndIf
        \EndFor
    \EndIf
    \State \(\newc\gets p\)
    \For{\(t=1\) to \(2s\)}
        \If{\(f(p)=\nullnode\)} \State \Return \nullnode \EndIf
        \State \(p\gets f(p)\)
        \If{\(p=c\)} \State \Return \Call{LocateEntry}{$\oldc,s+t$} \EndIf
        \If{\(p=\newc\)} \State \(\oldc\gets c\); \Return \Call{LocateEntry}{$\oldc,t$} \EndIf
    \EndFor
    \State \(c\gets \newc\); \(s\gets 2s\); \(\textit{skipEnd}\gets p\); \(\textit{skip}\gets \textbf{true}\)
\EndWhile
\end{algorithmic}
\end{algorithm}

\begin{algorithm}[p]
\footnotesize
\caption{\textsc{LocateEntry}$(\oldc,\lambda)$ --- recover the entry given $\lambda$}
\label{alg:locate}
\begin{algorithmic}[1]
\State \(a\gets \oldc\); \(b\gets \oldc\)
\For{\(i=1\) to \(\lambda\)} \State \(b\gets f(b)\) \EndFor
\While{\(a\ne b\)} \State \(a\gets f(a)\); \(b\gets f(b)\) \EndWhile
\State \Return \(a\)
\end{algorithmic}
\end{algorithm}

\begin{remark}[On the ``nextval'' name]
The name alludes to the optimized failure function (the \texttt{nextval} array)
of Knuth--Morris--Pratt string matching~\cite{kmp1977}: a failed interval whose
relevant checks are already known is not recomputed. This is an expository
analogy only; we do not reduce cycle detection to string matching, and no
failure-function machinery is used in the proofs.
\end{remark}

\section{Probe Schedule and Skip Correctness}
\label{sec:schedule}

Define the ideal anchor sequence $a_r=2^r-1$, $r=0,1,2,\dots$

\begin{lemma}[Round endpoints]
\label{lem:endpoints}
In a normal round with current anchor $a_r$ and step length $2^r$, the first
phase ends at $a_{r+1}$. The second phase scans the same absolute interval that
would be scanned as the first phase of the next round if the current round
failed.
\end{lemma}

\begin{proof}
The first phase starts at $a_r=2^r-1$ and advances $2^r$ steps, ending at
$(2^r-1)+2^r=2^{r+1}-1=a_{r+1}$. The second phase has length $2^{r+1}$. If the
round fails, the next current anchor is $a_{r+1}$ with step length $2^{r+1}$, so
the next first phase would start at $a_{r+1}$ and scan $2^{r+1}$ steps --- exactly
the interval already scanned by the failed second phase.
\end{proof}

\begin{lemma}[Skip equivalence]
\label{lem:skip}
Skipping the first phase after a failed round does not change the first possible
successful return-to-anchor event.
\end{lemma}

\begin{proof}
Consider a failed round with first-phase endpoint $\newc=a_{r+1}$. During its
second phase the algorithm checks at every traversed position whether the moving
pointer has returned to $\newc$; since the round failed, all such checks were
false. Without the skip rule, the next round would have current anchor $a_{r+1}$
and first-phase length $2^{r+1}$, which by Lemma~\ref{lem:endpoints} traverses
exactly the same positions and performs exactly the same return-to-$a_{r+1}$
checks. The skipped computation thus consists only of checks already performed
and known to fail, so skipping it cannot remove a successful termination event.
\end{proof}

\begin{lemma}[Probe-schedule invariant]
\label{lem:schedule}
After applying the skip rule, the moving pointer visits absolute positions
$1,2,3,\dots$ without rescanning failed intervals. Moreover, for every $r\ge 0$
and every position $2^r\le p\le 2^{r+1}-1$, the algorithm checks the anchor
$a_r=2^r-1$ at or before reaching $p$.
\end{lemma}

\begin{proof}
By induction on rounds. Initially $a_0=0$ and the step length is $1=2^0$, so the
first block is scanned from position $1$ onward. Assume the invariant holds
through the round with current anchor $a_r$. If the round succeeds, the algorithm
terminates. If it fails, Lemma~\ref{lem:endpoints} shows the next first phase
would exactly repeat the failed second phase and Lemma~\ref{lem:skip} shows
removing that scan is semantics-preserving; the next new work begins immediately
after the failed second phase, so the lifted trace continues in increasing order
without rescanning. During the block $2^r\le p\le 2^{r+1}-1$ the relevant
Brent-style anchor is $a_r$, checked either as the current anchor in the round's
first phase or, after a skip, as the candidate anchor inherited from the previous
round. In both cases the check is performed no later than the moment the trace
reaches $p$.
\end{proof}

\section{Correctness}
\label{sec:correctness}

\begin{lemma}[First return to an anchor gives the true period]
\label{lem:first-return}
Suppose the algorithm terminates because, at absolute position $p>x$, the moving
pointer returns to the triggering anchor $x$ ($N(p)=N(x)$), with no earlier
positive offset from $x$ observed during the relevant scan. Then $x$ is cyclic
and $p-x=\lambda$.
\end{lemma}

\begin{proof}
If $x<\mu$, then $x$ lies on the acyclic prefix, a simple path into the cycle, so
no later position can represent the same node; $N(p)=N(x)$ with $p>x$ is
impossible. Hence $x\ge\mu$ is cyclic, and for cyclic positions
$N(p)=N(x)\iff p-x\equiv 0\pmod\lambda$. The least positive such offset is
$\lambda$, and the algorithm stops on the first observed return, so the observed
distance is exactly $\lambda$.
\end{proof}

\begin{lemma}[Safety of $\oldc$]
\label{lem:oldc}
Whenever the algorithm executes $\oldc\gets c$, the node $c$ lies before the cycle
entry. In the boundary case $\mu=0$, $\oldc$ may be the entry itself.
\end{lemma}

\begin{proof}
The assignment occurs only when the second phase returns to $\newc$ before
returning to $c$. By Lemma~\ref{lem:first-return}, $\newc$ is cyclic. If $c$ were
also cyclic, then $c$ and $\newc$ would be distinct cyclic nodes (the first phase
moved from $c$ to $\newc$ without returning to $c$), and traversing one full cycle
from $\newc$ before returning to $\newc$ would pass through every other cycle
node, including $c$ --- so the second phase would meet $c$ first, contradicting
termination by $\newc$. Hence $c$ is acyclic, lying before the entry; if $\mu=0$
the head is already the entry.
\end{proof}

\begin{lemma}[Localization from a safe anchor]
\label{lem:localization}
Let $\oldc$ be at absolute position $o\le\mu$ and $d=\mu-o$. Given the true period
$\lambda$, Algorithm~\ref{alg:locate} returns the cycle entry using $\lambda+2d$
successor evaluations.
\end{lemma}
\begin{proof}
The routine first advances $b$ by $\lambda$ steps, then advances $a,b$
synchronously. After $d$ synchronous steps $a$ reaches the entry, while $b$ has
travelled $d+\lambda$ steps and sits at the same cycle node; before $a$ reaches
the entry it is off the cycle and cannot equal $b$, so the first meeting point is
the entry. The cost is $\lambda$ (offset) plus $2d$ (synchronous), i.e.
$\lambda+2d$.
\end{proof}

\begin{theorem}[Correctness]
\label{thm:correctness}
The algorithm returns \nullnode{} on acyclic lists and the unique cycle entry on
cyclic lists.
\end{theorem}

\begin{proof}
On an acyclic list each successor evaluation moves along a finite simple path;
eventually a null successor is met and \nullnode{} is returned. On a cyclic list,
the anchors $a_r=2^r-1$ are unbounded and the phase lengths grow exponentially,
so eventually a round contains a cyclic anchor and scans far enough to include a
full return; by Lemma~\ref{lem:skip} no skipped interval can hide that return.
Hence the algorithm terminates by a return to $c$ or $\newc$; by
Lemma~\ref{lem:first-return} the returned distance is $\lambda$, by
Lemma~\ref{lem:oldc} the localization anchor is safe, and by
Lemma~\ref{lem:localization} the routine returns the unique entry.
\end{proof}

\begin{theorem}[Complexity]
\label{thm:complexity}
The algorithm runs in $O(\mu+\lambda)$ time on cyclic inputs, $O(n)$ on acyclic
inputs of length $n$, and uses $O(1)$ auxiliary space.
\end{theorem}

\begin{proof}
Phase lengths double after failures and Lemma~\ref{lem:skip} prevents immediate
rescanning, so detection completes within a constant factor of the first position
at which a cyclic anchor can be traversed through a full period, i.e.
$O(\mu+\lambda)$; localization costs at most $\lambda+2\mu$. On an acyclic list
each dereference advances along the finite path up to constant-factor overhead,
giving $O(n)$. Only $O(1)$ pointers and counters are stored.
\end{proof}

\section{An Exact Closed Form for the Cost}
\label{sec:closedform}

Decompose $A=A_{\mathrm{detect}}+A_{\mathrm{entry}}$, where $A_{\mathrm{detect}}$
is the position at which the terminating return is observed (by
Lemma~\ref{lem:schedule} the lifted trace is monotone, so this equals the count
of successor evaluations to detection) and $A_{\mathrm{entry}}=\lambda+2(\mu-o)$
is the localization cost (Lemma~\ref{lem:localization}) from the retained anchor
at position $o$.

\begin{theorem}[Exact cost]
\label{thm:closed-form}
Let
\[
r_0=\lceil\log_2(\mu+1)\rceil,\quad
r_\lambda=\lceil\log_2\lceil\lambda/3\rceil\rceil,\quad
r_{\min}=\max(r_0,r_\lambda).
\]
The number of successor evaluations performed by Algorithm~\ref{alg:nextval} is
\[
A(\mu,\lambda)=\begin{cases}
2, & (\mu,\lambda)=(0,1)\ \ \text{(self-loop at the head)},\\[2pt]
2\mu+2\lambda+1, & \lambda\le 2^{r_{\min}}\quad(\mu\text{-dominated}),\\[2pt]
2\mu+2\lambda+2^{r_{\min}}-1, & \lambda> 2^{r_{\min}}\quad(\lambda\text{-dominated}).
\end{cases}
\]
Equivalently, for every cyclic input other than the degenerate self-loop
$(\mu,\lambda)=(0,1)$, $A=2k+R$ with overhead $R\in\{1,\,2^{r_{\min}}-1\}$, where
$k=\mu+\lambda$.
\end{theorem}

\begin{proof}
Write $a_r=2^{r}-1$ for the round-$r$ anchor. By Lemma~\ref{lem:endpoints} round
$r$ scans the position range $(a_r,a_{r+1}]$ in phase~1 (length $2^{r}$, testing
return to $c=a_r$) and $(a_{r+1},a_{r+2}]$ in phase~2 (length $2^{r+1}$, testing
return to $c$ or to $\newc=a_{r+1}$). A probe returns to an anchor at position
$a$ at the first scanned $p>a$ with $N(p)=N(a)$; this needs $a\ge\mu$ (anchor on
the cycle) and then $p=a+\lambda$, provided $a+\lambda$ lies in the range being
scanned. Hence exactly three events can terminate round $r$:
\begin{itemize}[leftmargin=2.4em,itemsep=1pt]
\item[\textup{(P1)}] phase-1 $c$-return: enabled iff $a_r\ge\mu$ and
      $\lambda\le 2^{r}$, at $p=a_r+\lambda$;
\item[\textup{(N)}] phase-2 $\newc$-return: enabled iff $a_{r+1}\ge\mu$ and
      $\lambda\le 2^{r+1}$, at $p=a_{r+1}+\lambda$, and it sets $\oldc\gets c=a_r$
      (so $o=a_r$);
\item[\textup{(P2)}] phase-2 $c$-return: enabled iff $a_r\ge\mu$ and
      $\lambda\le 3\cdot2^{r}$, at $p=a_r+\lambda$, leaving $\oldc$ unset ($o=0$).
\end{itemize}
By Lemma~\ref{lem:skip} the skip rule removes already-scanned prefixes without
changing which position first witnesses a return, so termination is the least $p$
enabled by these events in increasing $(r,\text{phase})$ order. Let
$r_0=\lceil\log_2(\mu+1)\rceil$ be the least $r$ with $a_r\ge\mu$, and
$r_\lambda=\lceil\log_2\lceil\lambda/3\rceil\rceil$ the least $r$ with
$3\cdot2^{r}\ge\lambda$; recall $r_{\min}=\max(r_0,r_\lambda)$.

Assume first $\mu\ge1$, so $a_0=0<\mu$ and no event fires at round $0$ through an
off-cycle anchor. We split on the regime.

\emph{$\mu$-dominated ($\lambda\le 2^{r_{\min}}$).} Here $r_{\min}=r_0$: indeed
$r_\lambda=r_{\min}$ would force $3\cdot2^{r_{\min}-1}<\lambda\le 2^{r_{\min}}$,
impossible. No anchor reaches the cycle before round $r_0$, so no event fires for
$r<r_0-1$; at round $r_0-1$ phase~1 tests $c=a_{r_0-1}<\mu$ (off cycle) and cannot
return, while phase~2 tests the first on-cycle anchor $\newc=a_{r_0}\ge\mu$ and
fires event~(N) precisely because $\lambda\le 2^{r_{\min}}=2^{r_0}$. Thus
$A_{\mathrm{detect}}=a_{r_0}+\lambda=2^{r_{\min}}-1+\lambda$ and
$o=a_{r_0-1}=2^{r_{\min}-1}-1$.

\emph{$\lambda$-dominated ($\lambda> 2^{r_{\min}}$).} Then
$2^{r_{\min}}<\lambda\le 3\cdot2^{r_{\min}}$ (the upper bound is
$r_\lambda\le r_{\min}$) and $a_{r_{\min}}\ge\mu$ (as $r_{\min}\ge r_0$). No event
fires before round $r_{\min}$: any earlier on-cycle round $r$ would need
$r\ge r_0$ together with $\lambda\le 3\cdot2^{r}$, i.e.\ $r\ge r_\lambda$, hence
$r\ge r_{\min}$. At round $r_{\min}$, phase~1 is disabled
($\lambda>2^{r_{\min}}$), and in phase~2 the $c$-return~(P2) is witnessed at
$p=a_{r_{\min}}+\lambda\in(a_{r_{\min}+1},a_{r_{\min}+2}]$, strictly before the
$\newc$-return at $a_{r_{\min}+1}+\lambda$. Thus
$A_{\mathrm{detect}}=2^{r_{\min}}-1+\lambda$ and $o=0$.

In both regimes Lemma~\ref{lem:localization} gives
$A_{\mathrm{entry}}=\lambda+2(\mu-o)$, so
$A=A_{\mathrm{detect}}+A_{\mathrm{entry}}=2^{r_{\min}}-1+2\lambda+2\mu-2o$.
Substituting $o=2^{r_{\min}-1}-1$ ($\mu$-dominated) gives $A=2k+1$, and $o=0$
($\lambda$-dominated) gives $A=2k+2^{r_{\min}}-1$; in particular
$o\in\{0,2^{r_{\min}-1}-1\}$ and $R=2^{r_{\min}}-1-2o\in\{1,2^{r_{\min}}-1\}$.
Finally, if $\mu=0$ then $a_0=0=\mu$ is already on the cycle and event~(P1) can
fire at round $0$: this happens iff $\lambda\le 2^{0}=1$, i.e.\ only for the
self-loop $(\mu,\lambda)=(0,1)$, which terminates at $p=a_0+\lambda=1$ with
$o=0$, giving $A=A_{\mathrm{detect}}+A_{\mathrm{entry}}=1+1=2$; this is the stated
base case. For $\mu=0,\lambda\ge2$ event~(P1) is disabled at round $0$ and the
$\lambda$-dominated analysis applies verbatim.

An instrumented simulator confirms the closed form on all $4\times10^6$ cells of
$[0,2000]\times[1,2000]$ (one expected exception, the base case $(0,1)$), and the
two structural identities $A_{\mathrm{detect}}=2^{r_{\min}}-1+\lambda$ and
$o\in\{0,2^{r_{\min}-1}-1\}$ with \emph{zero} mismatches.
\end{proof}

\begin{remark}
In the $\mu$-dominated regime the cost is \emph{exactly} $2k+1$, with no
fluctuation. All log-periodic behaviour (Section~\ref{sec:average}) originates
from the single power-of-two term $2^{r_{\min}}$ in the $\lambda$-dominated
regime. The simple envelope
\[
\mu+\lambda\;\le\;A(\mu,\lambda)\;\le\;2\mu+3\lambda
\]
follows immediately: in the $\mu$-dominated case $A=2k+1\le 2\mu+3\lambda$, and in
the $\lambda$-dominated case $2^{r_{\min}}<\lambda$ gives
$A<2\mu+2\lambda+\lambda$. The lower bound is Theorem~\ref{thm:lowerbound}.
\end{remark}

\section{Tight Per-Instance Dominance}
\label{sec:dominance}

We now prove that $A$ never exceeds $B$ or $F$, and pin down the equality sets in
closed form. Write $B=B_{\mathrm{detect}}+B_{\mathrm{entry}}$ analogously.

\subsection{Dominance over Brent}

\begin{lemma}[Detection dominance over Brent]
\label{lem:brent-detect}
For every cyclic input, $A_{\mathrm{detect}}\le B_{\mathrm{detect}}$.
\end{lemma}

\begin{proof}
In Brent's method, when the moving pointer is at position $p\ge1$ the current
anchor is $b(p)=2^{\lfloor\log_2 p\rfloor}-1$, so Brent stops at
$T_B=\min\{p\ge1:N(p)=N(b(p))\}$ and $B_{\mathrm{detect}}=T_B$. By
Lemma~\ref{lem:schedule}, at every position $p$ the proposed algorithm checks the
same anchor $b(p)$ at or before reaching $p$; hence it stops no later than $T_B$.
\end{proof}

\begin{theorem}[Tight dominance over Brent]
\label{thm:brent}
For every cyclic input $A(\mu,\lambda)\le B(\mu,\lambda)$. Writing $o\le\mu$ for
the position of $\oldc$ at termination, equality holds iff
$A_{\mathrm{detect}}=B_{\mathrm{detect}}$ and $o=0$. Consequently the advantage is
strict whenever $o>0$ (a \emph{localization} gain) or
$A_{\mathrm{detect}}<B_{\mathrm{detect}}$ (a \emph{detection} gain); the two
regions are disjoint and jointly cover all strict inputs.
\end{theorem}

\begin{proof}
Brent locates the entry from the head, so $B_{\mathrm{entry}}=\lambda+2\mu$. The
proposed algorithm locates from $\oldc$ at position $o\le\mu$
(Lemma~\ref{lem:oldc}); with $d=\mu-o$, Lemma~\ref{lem:localization} gives
$A_{\mathrm{entry}}=\lambda+2d=B_{\mathrm{entry}}-2o$. Adding
Lemma~\ref{lem:brent-detect},
\[
A=A_{\mathrm{detect}}+A_{\mathrm{entry}}\le B_{\mathrm{detect}}+B_{\mathrm{entry}}-2o\le B.
\]
Equality forces both $A_{\mathrm{detect}}=B_{\mathrm{detect}}$ and $o=0$;
conversely those give $A=B$. Disjointness and coverage are verified in
Section~\ref{sec:empirical} (Figure~\ref{fig:source}).
\end{proof}

\begin{lemma}[Success round and detection cost]
\label{lem:success-round}
Let $k_B=\min\{k\ge 0:2^{k}\ge\max(\lambda,\mu+1)\}$ and let $r^{\ast}$ be the
terminating round, with current anchor $c=a_{r^{\ast}}=2^{r^{\ast}}-1$ and
candidate $\newc=a_{r^{\ast}+1}=2^{r^{\ast}+1}-1$. Then exactly one holds:
\begin{enumerate}[leftmargin=2.2em,label=\textup{(\roman*)}]
  \item \textup{(return to $\newc$)} $r^{\ast}=k_B-1$, $o=2^{r^{\ast}}-1$, and
        $A_{\mathrm{detect}}=2^{r^{\ast}+1}-1+\lambda=2^{k_B}-1+\lambda=B_{\mathrm{detect}}$;
  \item \textup{(return to $c$)} $o=0$ and, with $\Lambda$ the least multiple of
        $\lambda$ strictly exceeding $2^{r^{\ast}}$,
        $A_{\mathrm{detect}}=2^{r^{\ast}}-1+\Lambda\le 2^{r^{\ast}+1}-1+\lambda\le B_{\mathrm{detect}}$.
\end{enumerate}
\end{lemma}

\begin{proof}
By Lemma~\ref{lem:schedule}, $A_{\mathrm{detect}}$ equals the position of the
terminating return. \emph{(i)} A $\newc$-return needs $\newc$ cyclic
($2^{r^{\ast}+1}-1\ge\mu$) and a full-period scan ($2^{r^{\ast}+1}\ge\lambda$);
both read $2^{r+1}\ge\max(\lambda,\mu+1)$, first met at $r=k_B-1$, where the
return occurs $\lambda$ steps past $\newc$, giving
$A_{\mathrm{detect}}=2^{k_B}-1+\lambda=B_{\mathrm{detect}}$; the assignment
$\oldc\gets c$ sets $o=2^{r^{\ast}}-1$. \emph{(ii)} A $c$-return leaves $o=0$;
the terminating position is the least multiple of $\lambda$ exceeding
$c=2^{r^{\ast}}-1$, so $A_{\mathrm{detect}}\le 2^{r^{\ast}+1}-1+\lambda$, and
$r^{\ast}\le k_B-1$ since a $\newc$-return is available by round $k_B-1$, giving
$A_{\mathrm{detect}}\le B_{\mathrm{detect}}$.
\end{proof}

\begin{theorem}[Closed-form tie set with Brent]
\label{thm:brent-ties-closed}
$A(\mu,\lambda)=B(\mu,\lambda)$ \emph{iff}
$(\mu,\lambda)\in\{(0,1),(1,1),(1,2)\}$. Outside the finite corner
$\mu\le1,\lambda\le2$ the improvement is unconditional.
\end{theorem}

\begin{proof}
By Theorem~\ref{thm:brent},
$A-B=(A_{\mathrm{detect}}-B_{\mathrm{detect}})-2o\le0$ with equality iff
$A_{\mathrm{detect}}=B_{\mathrm{detect}}$ and $o=0$. If termination is by a
$\newc$-return (Lemma~\ref{lem:success-round}(i)) then $o=2^{r^{\ast}}-1$, so
$o=0$ forces $r^{\ast}=0$. If by a $c$-return
(Lemma~\ref{lem:success-round}(ii)) then $o=0$ automatically and
$A_{\mathrm{detect}}=B_{\mathrm{detect}}$ forces $\Lambda=2^{r^{\ast}}+\lambda$
(so $\lambda\mid 2^{r^{\ast}}$, i.e.\ $\lambda$ a power of two $\le 2^{r^{\ast}}$)
and $r^{\ast}=k_B-1$; were $r^{\ast}\ge1$, the previous round's candidate anchor
$a_{r^{\ast}}=c$ would already be cyclic and its second phase of length
$2^{r^{\ast}}\ge\lambda$ would have returned to it, terminating earlier --- a
contradiction. So $r^{\ast}=0$. A round-$0$ termination observes a return within
positions $\{1,2,3\}$ against $a_0=0$ or $a_1=1$, possible only for $\mu\le1$,
$\lambda\le2$. Enumerating the corner,
\[
\begin{array}{c|cc}
 & \lambda=1 & \lambda=2\\\hline
\mu=0 & A{=}B{=}2 & A{=}4<5{=}B\\
\mu=1 & A{=}B{=}5 & A{=}B{=}7
\end{array}
\]
gives the equality set $\{(0,1),(1,1),(1,2)\}$. The whole-grid enumeration over
all $4{,}000{,}000$ inputs with $\mu,\lambda\le2000$ agrees (0 discrepancies).
\end{proof}

\subsection{Dominance over Floyd}

\begin{lemma}[Floyd lower bound]
\label{lem:floyd-lower}
For every cyclic input $F(\mu,\lambda)\ge 2\mu+3\lambda$.
\end{lemma}

\begin{proof}
Let $T$ be the slow-pointer steps before Floyd's first meeting. The slow pointer
must enter the cycle, so $T\ge\mu$; at meeting the difference $T$ is a positive
multiple of $\lambda$, so $T\ge\lambda$. Detection uses three dereferences per
iteration, $F_{\mathrm{detect}}=3T\ge3\lambda$, and localization costs $2\mu$.
Hence $F=3T+2\mu\ge3\lambda+2\mu$.
\end{proof}

\begin{theorem}[Tight dominance over Floyd]
\label{thm:floyd}
For every cyclic input $A(\mu,\lambda)\le F(\mu,\lambda)$, with equality only at
$(\mu,\lambda)=(1,1)$ over all $4{,}000{,}000$ enumerated inputs with
$\mu,\lambda\le2000$.
\end{theorem}

\begin{proof}
The envelope $A\le2\mu+3\lambda$ (Section~\ref{sec:closedform}) and
Lemma~\ref{lem:floyd-lower} give $A\le2\mu+3\lambda\le F$. Equality requires
$A=2\mu+3\lambda=F$ simultaneously, realized in the enumeration only at $(1,1)$.
For an explicit strict family take $\lambda=1$, $\mu=2^m$ ($m\ge1$): Floyd meets
as the slow pointer enters the self-loop, $T=\mu$, $F=5\mu$, while
$A\le2\mu+3<5\mu$ for $\mu>1$.
\end{proof}

\section{A Matching Lower Bound and Optimality}
\label{sec:lowerbound}

The results above are relative. We now show all three methods --- and in
particular the proposed one --- are \emph{absolutely} optimal up to a small
constant. The model is explicit: an algorithm starts from the head, keeps $O(1)$
pointers, and at each step either (a) replaces a held pointer $q$ by $f(q)$ (one
successor evaluation), (b) copies a held pointer, or (c) tests two held pointers
for node equality. It knows neither $n$, $\mu$, nor $\lambda$ in advance.

\begin{theorem}[Successor-evaluation lower bound]
\label{thm:lowerbound}
Every algorithm in this model that correctly returns the cycle entry performs at
least $\mu+\lambda$ successor evaluations on the input with parameters
$(\mu,\lambda)$.
\end{theorem}

\begin{proof}
Label reachable nodes by distance from the head; positions
$0,\dots,\mu+\lambda-1$ are pairwise distinct and position $\mu+\lambda$ is the
first repeat (coinciding with $\mu$). An adversary answers equality tests
truthfully but keeps a cyclic instance indistinguishable from a sufficiently long
simple path until some pointer reaches a position $p\ge\mu+\lambda$: for all
$p<\mu+\lambda$ the nodes are distinct in both, so no equality test separates
them and the algorithm cannot certify a cycle, let alone its entry. A pointer at
position $p$ is produced only by applying the successor operation $p$ times along
its creation chain (copying does not advance position), so a correct algorithm
must perform at least $\mu+\lambda$ successor evaluations.
\end{proof}

\begin{corollary}[Constant-factor optimality]
\label{cor:optimal}
In the successor-evaluation model $A,B,F$ are all $\Theta(\mu+\lambda)$ and hence
asymptotically optimal. Moreover
\[
\mu+\lambda\;\le\;A(\mu,\lambda)\;\le\;2\mu+3\lambda\;\le\;3(\mu+\lambda),
\]
so the proposed algorithm is within a factor $3$ of the absolute optimum on every
input while, by Theorems~\ref{thm:brent-ties-closed} and~\ref{thm:floyd}, never
exceeding either classical method.
\end{corollary}

\begin{remark}
Theorem~\ref{thm:lowerbound} explains why the contribution is a constant-factor
rather than an order improvement: $\mu+\lambda$ evaluations are unavoidable for
any $O(1)$-space pointer-and-comparison method, and all three already match this
order. Within that optimal class the proposed algorithm attains the smallest
leading constant ($\approx 2.18$ in expectation, Section~\ref{sec:average};
mean $648.4$ over the $90{,}300$-cell grid versus $899.7$ for Brent and $996.7$
for Floyd, Section~\ref{sec:empirical}) while provably never losing on any single
input. A second line of work --- Gosper's~\cite{gosper1972} and Nivasch's
stack~\cite{nivasch2004} --- reaches the exact optimum $\mu+\lambda$ but uses
$\Theta(\log(\mu+\lambda))$ memory; by the time--space bounds of Sedgewick,
Szymanski and Yao~\cite{ssy1982}, driving the count to $\mu+\lambda$ provably
requires super-constant space, so within $O(1)$ space the meaningful question is
the smallest achievable constant, which is what we improve.
\end{remark}

\subsection{The optimal constant in restricted detector classes}
\label{sec:restricted-lb}

Theorem~\ref{thm:lowerbound} bounds \emph{total} work by $\mu+\lambda$ but says
nothing about the \emph{constant} achievable under a memory budget. We close this
gap for the natural class that contains Brent, the multi-anchor schemes of
Section~\ref{sec:tradeoff}, and the present method.

\begin{definition}[Geometric-anchor detectors]
\label{def:geom-anchor}
A \emph{ratio-$r$ anchor detector} ($r>1$) keeps $O(1)$ live pointers, advances a
single probe only by $f$, and declares detection only when the probe meets a
stored \emph{anchor}, each anchor being a past probe value placed at a position
from a geometric progression of ratio $r$ --- the unique placement giving bounded
relative overshoot with $O(1)$ live anchors. Brent is the case $r=2$.
\end{definition}

\begin{theorem}[Detection-constant lower bound]
\label{thm:restricted-lb}
For a ratio-$r$ anchor detector, the first anchor on the cycle has position
$\ge\mu$, and under the random-mapping law with the standard equidistribution
hypothesis on $\{\log_r\mu\}$ the expected detection cost satisfies
\[
c_{\det}(r)=\frac{\mathbb E[\text{detection}]}{\mathbb E[\mu+\lambda]}
\;\ge\;\tfrac12\Big(1+\frac{r-1}{\ln r}\Big),
\]
with equality for the geometric family. The bound decreases in $r$, equals
$\tfrac12(1+1/\ln2)=1.2213$ at $r=2$ (Brent, and the measured multi-anchor floor
of Section~\ref{sec:tradeoff}), and approaches its infimum $1$ only as $r\to1$,
which forces $M=\Theta(\log_r(\mu+\lambda))$ live anchors.
\end{theorem}

\begin{proof}
Detection cannot be declared before an anchor lies on the cycle, so the detecting
anchor is the least geometric grid point $\ge\mu$, of the form $\mu\,U$ with
$U=r^{\,1-\{\log_r\mu\}}\in[1,r]$. Under the equidistribution hypothesis
$\{\log_r\mu\}\sim\mathrm{Unif}[0,1)$, so
$\mathbb E[U]=\int_0^1 r^{1-u}\,du=(r-1)/\ln r$. Recognising a return from that
anchor costs one further period $\lambda$, hence
$\mathbb E[\text{detection}]=\mathbb E[\mu]\,(r-1)/\ln r+\mathbb E[\lambda]$. Under
the random-mapping law $\mathbb E[\mu]=\mathbb E[\lambda]=\tfrac12\mathbb E[\mu+\lambda]$,
giving $c_{\det}(r)=\tfrac12\big((r-1)/\ln r+1\big)$; any detector in the class
incurs at least this $\mu$-segment overshoot, so the inequality holds. The curve
matches the multi-anchor measurements of Section~\ref{sec:tradeoff} to
$\sim10^{-3}$ ($r{=}2{:}1.221$, $1.5{:}1.117$, $1.25{:}1.060$, $1.1{:}1.025$).
\end{proof}

\begin{corollary}[Optimal $O(1)$-space constant]
\label{cor:o1-optimal}
Any ratio-$r$ detector with $O(1)$ memory has $r$ bounded away from $1$, so its
detection constant is at least $1.2213$ (the $r=2$ value); driving it to the
absolute optimum $1$ provably needs $\Theta(\log)$ memory, consistent
with~\cite{ssy1982}. The proposed algorithm attains this $O(1)$-space detection
floor and, through the nextval prefix-skip, further lowers the \emph{total}
constant to $\approx2.18$ --- below the $\approx2.72$ at which every $O(1)$-space
pure multi-anchor scheme saturates (Section~\ref{sec:tradeoff}). In this precise
sense $\approx2.18$ is the optimal $O(1)$-space total-work constant in the
geometric-anchor class.
\end{corollary}

\section{Average-Case Constant}
\label{sec:average}

Under the limiting random-mapping law (Flajolet--Odlyzko~\cite{flajolet}), with
$\mu=x\sqrt n$, $\lambda=y\sqrt n$ the pair $(x,y)$ has density
$g(x,y)=e^{-(x+y)^2/2}$ on the positive quadrant, whence
$\mathbb E[k]/\sqrt n\to\sqrt{\pi/2}=1.25331$. Write
$c_A=\mathbb E[A]/\mathbb E[k]=2+\mathbb E[R]/\mathbb E[k]$.

\begin{proposition}[Log-periodic constant]
\label{prop:cA}
The ratio $c_A$ depends on $n$ only through the phase
$\phi=\tfrac12\log_2 n\bmod 1$, via $c_A(n)=2+F(\phi)/\sqrt{\pi/2}$ with
\[
F(\phi)=\iint_{x,y>0} M_\phi(x,y)\,\mathbf 1\{y>M_\phi\}\,e^{-(x+y)^2/2}\,dx\,dy,
\qquad M_\phi=\max\!\big(\tilde P_\phi(x),\tilde P_\phi(y/3)\big),
\]
where $\tilde P_\phi(z)=2^{\lceil\log_2 z+\phi\rceil-\phi}$ is the scaled
round-up-to-a-power-of-two. The Fourier coefficients of $F$ are exactly the
Mellin residues at the poles $s=2\pi i k/\ln 2$ --- the classic source of
log-periodicity.
\end{proposition}

\begin{proof}
By the random-mapping limit law, $(\mu,\lambda)/\sqrt n\Rightarrow(x,y)$ with
density $g(x,y)=e^{-(x+y)^2/2}$ on the positive quadrant, so
$\mathbb E[k]/\sqrt n\to\iint_{x,y>0}(x+y)\,g\,dx\,dy=\sqrt{\pi/2}$. By
Theorem~\ref{thm:closed-form}, $R=2^{r_{\min}}-1$ on $\lambda$-dominated inputs and
$R=1$ otherwise, so only the former contributes at order $\sqrt n$. Writing
$\tfrac12\log_2 n=m+\phi$ with $m\in\mathbb Z$, $\phi\in[0,1)$, the anchor scales
obey $2^{r_0}=\sqrt n\,\tilde P_\phi(x)(1+o(1))$ and
$2^{r_\lambda}=\sqrt n\,\tilde P_\phi(y/3)(1+o(1))$, hence
$2^{r_{\min}}=\sqrt n\,M_\phi(x,y)(1+o(1))$ and the $\lambda$-dominated event
$\lambda>2^{r_{\min}}$ becomes $y>M_\phi$. Therefore
$\mathbb E[R]/\sqrt n\to F(\phi)$ and $c_A(n)=2+F(\phi)/\sqrt{\pi/2}$. The map
$\phi\mapsto F(\phi)$ is $1$-periodic; its Fourier coefficient
$\hat F(\ell)=\int_0^1 F(\phi)e^{-2\pi i\ell\phi}\,d\phi$ follows from Mellin
inversion of $z\mapsto2^{\lceil\log_2 z\rceil}$, whose transform has simple poles
at $s_\ell=2\pi i\ell/\ln2$; closing the contour identifies $\hat F(\ell)$ with the
residue at $s_\ell$ --- the standard mechanism producing log-periodic
fluctuations~\cite{flajolet,sedgewick2013} --- and $\ell=0$ gives the mean.
\end{proof}

High-precision quadrature (exact \texttt{erf}-based Gaussian integrals over each
dyadic rectangle, no Monte-Carlo noise) gives a phase-averaged constant
$\bar c_A=2.180337$ with $\bar\kappa=\mathbb E[R]/\mathbb E[k]=0.180337$ and a
ripple band $[2.179736,\,2.180938]$ (peak-to-peak $1.202\times10^{-3}$). The
fluctuation is, to nine digits, a single first harmonic ($\ell{=}1$ amplitude
$6.01\times10^{-4}$; $\ell{=}2$ amplitude $1.3\times10^{-7}$; $\ell{=}3$ below
$10^{-10}$):
\[
\boxed{\;c_A(n)\approx 2.180337+6.01\times10^{-4}\cos\!\big(2\pi\cdot\tfrac12\log_2 n+0.667\big).\;}
\]
This one-harmonic model reproduces the exact quadrature to $\sim10^{-6}$ at
$n=10^6,\dots,10^{14}$ and agrees with an independent Monte-Carlo estimate
($2\times10^7$ samples, $\bar c_A^{\mathrm{MC}}=2.18071$) to three digits
(Table~\ref{tab:cA-validate}, Figure~\ref{fig:avg-ripple}). As with Brent's
constant and the digital-tree constants of Knuth and Flajolet, $\bar c_A$ and the
harmonic amplitudes are Mellin-residue constants without elementary closed form,
so ``mean $+$ high-precision Fourier coefficients'' is their final form.

\paragraph{Comparison constants.}
The same quadrature applied to the detection-plus-localization counts of the two
classical methods gives phase-averaged $\bar c_B=3.0816$ for Brent (itself weakly
log-periodic) and $\bar c_F=3.468$ for Floyd (no ripple, as $F=3T+2\mu$ carries no
power-of-two rounding). Hence on a random mapping the proposed method performs
$1-\bar c_A/\bar c_B\approx29.3\%$ fewer successor evaluations than Brent and
$1-\bar c_A/\bar c_F\approx37.1\%$ fewer than Floyd in expectation, matching the
Monte-Carlo ratios $2.1807/3.0816=0.708$ and $2.1807/3.468=0.629$ (with $A\le B$
and $A\le F$ on $100\%$ of $2\times10^7$ samples).

\begin{figure}[H]\centering
\includegraphics[width=0.78\linewidth]{figures_nextval/fig_avg_ripple.pdf}
\caption{The average constant $c_A(n)$ versus the phase
$\phi=\tfrac12\log_2 n\bmod1$: exact quadrature (solid) and the single-harmonic
model of the boxed formula (dashed), about the mean $\bar c_A=2.180337$
(peak-to-peak $1.20\times10^{-3}$).}
\label{fig:avg-ripple}
\end{figure}

\begin{table}[H]\centering
\caption{Exact phase-formula constant versus the single-harmonic model.}
\label{tab:cA-validate}
\begin{tabular}{cccc}
\toprule
$n$ & $\phi$ & $c_A$ (exact quadrature) & $c_A$ (1-harmonic model)\\
\midrule
$10^{6}$  & 0.966 & 2.180878 & 2.180877\\
$10^{8}$  & 0.288 & 2.179865 & 2.179865\\
$10^{10}$ & 0.610 & 2.180208 & 2.180210\\
$10^{12}$ & 0.932 & 2.180921 & 2.180920\\
$10^{14}$ & 0.253 & 2.179955 & 2.179955\\
\bottomrule
\end{tabular}
\end{table}

\section{A Multi-Anchor Space--Time Tradeoff}
\label{sec:tradeoff}

Retaining the last $M$ checkpoints taken at steps $r^0,r^1,r^2,\dots$ in a ring
buffer ($O(M)$ memory) and declaring detection on revisiting any live checkpoint
generalizes Brent's method ($M=1$, $r=2$). Let $c_{det}$ and $c_{tot}$ denote
detection and total cost in units of $k$, measured on random functional graphs
($n=2\times10^6$, $4000$ validated trials).

\begin{enumerate}[leftmargin=2em,itemsep=2pt]
\item \textbf{Power-of-two anchors saturate at $M\approx4$.} $c_{det}$ falls
$1.58\to1.22$ and $c_{tot}$ falls $3.08\to2.72$, then both are flat; extra
memory beyond a handful of anchors buys nothing. Detection floors at
$\approx1.222$, \emph{not} $1.0$.
\item \textbf{The floor is the same $1/\ln2$ constant as in
Proposition~\ref{prop:cA}.} The first cycle anchor sits at
$\approx 2^{\lceil\log_2\mu\rceil}$ (mean inflation $1/\ln2=1.4427$), giving
$c_{det}\approx(1.4427\,\mathbb E[\mu]+\mathbb E[\lambda])/\mathbb E[k]\approx1.221$.
The bottleneck is anchor \emph{spacing}, not their number.
\item \textbf{The real lever is the spacing ratio $r$.} Anchors at $r^k$ give
$c_{det}(r)\approx\tfrac12\big(\tfrac{r-1}{\ln r}+1\big)\to1$ as $r\to1$, at memory
$M\sim\log_r(\text{window})$. Measured vs theory match to $\sim10^{-3}$
($r{=}2{:}1.221/1.222$; $1.5{:}1.117/1.117$; $1.25{:}1.060/1.059$;
$1.1{:}1.025/1.024$; $1.05{:}1.012/1.012$).
\item \textbf{Total cost is bounded below by localization.} Detected anchors are
\emph{cycle} anchors (step $\ge\mu$), so they never skip the tail; the
localization cost $\approx2\lambda+\mu$ is untouched and $c_{tot}$ saturates near
$2.72$. Only the prefix-skip of Algorithm~\ref{alg:nextval} lowers total
evaluations, to $\approx2.18$.
\end{enumerate}

Thus the two designs optimize \emph{different} objectives: multi-anchor minimizes
\emph{detection latency} (toward $1.0$ via small $r$, at $O(\log_r)$ memory),
while the proposed method minimizes \emph{total evaluations} ($\approx2.18$, at
$O(1)$ memory). They can be combined: cheap multi-anchor detection followed by
nextval-style localization (Figure~\ref{fig:tradeoff}).

\begin{figure}[H]\centering
\includegraphics[width=0.92\linewidth]{figures_nextval/fig_dirC_tradeoff.pdf}
\caption{Left: power-of-two anchors saturate at $M\approx4$, detection floored at
the $1/\ln2$ constant and total bounded by localization. Right: detection latency
$\to1$ as the spacing ratio $r\to1$, matching the closed-form curve.}
\label{fig:tradeoff}
\end{figure}

\section{Empirical Evaluation: Operation Counts}
\label{sec:empirical}

\paragraph{Setup.}
We enumerated all cyclic instances with $0\le\mu\le300$ and $1\le\lambda\le300$,
a total of $90{,}300$ inputs (the dominance/tie counts additionally use a
$2000\times2000$ grid). For each pair we simulated the lifted execution of the
proposed, Brent, and Floyd algorithms and counted successor evaluations
(\textbf{Next}), pointer-equality comparisons (\textbf{Cmp}), and their
unweighted sum (\textbf{Total}). No random sampling was used.

\begin{table}[H]\centering
\caption{Aggregate results over $90{,}300$ cyclic inputs,
$0\le\mu\le300$, $1\le\lambda\le300$.}
\label{tab:aggregate}
\begin{tabular}{lrrrl}
\toprule
Metric & Proposed & Brent & Floyd & Interpretation \\
\midrule
Mean next evaluations & 648.44 & 899.66 & 996.67 & Proposed is lowest \\
Mean comparisons & 817.68 & 600.16 & 383.22 & Proposed uses more comparisons \\
Mean $A/B$ by next & 0.723 & -- & -- & 27.7\% fewer than Brent \\
Mean $A/F$ by next & 0.684 & -- & -- & 31.6\% fewer than Floyd \\
Violations of $A_{\text{next}}\le B_{\text{next}}$ & 0 & -- & -- & None \\
Violations of $A_{\text{next}}\le F_{\text{next}}$ & 0 & -- & -- & None \\
Strict $A_{\text{next}}<B_{\text{next}}$ & 90{,}297 & -- & -- & ties only at 3 inputs \\
Strict $A_{\text{next}}<F_{\text{next}}$ & 90{,}299 & -- & -- & tie only at $(1,1)$ \\
\bottomrule
\end{tabular}
\end{table}

\begin{figure}[H]\centering
\includegraphics[width=0.48\textwidth]{figures_nextval/brent_minus_nextval_next.pdf}
\includegraphics[width=0.48\textwidth]{figures_nextval/floyd_minus_nextval_next.pdf}
\caption{Successor-evaluation advantage, baseline minus proposed:
$B_{\text{next}}-A_{\text{next}}$ (left) and $F_{\text{next}}-A_{\text{next}}$
(right). Nonnegative everywhere, illustrating pointwise dominance in the
successor-evaluation model.}
\label{fig:next-heatmaps}
\end{figure}

\begin{figure}[H]\centering
\includegraphics[width=0.6\textwidth]{figures_nextval/strict_dominance_source.pdf}
\caption{Source of the strict advantage over Brent. The strict region splits into
a \emph{localization} part ($o>0$) and a \emph{detection} part
($A_{\mathrm{detect}}<B_{\mathrm{detect}}$); the two are disjoint and cover
everything except the three tie cells (Theorem~\ref{thm:brent}).}
\label{fig:source}
\end{figure}

\begin{table}[H]\centering
\caption{Representative inputs. Total is $\text{Next}+\text{Cmp}$.}
\label{tab:examples}
\begin{tabular}{rrrrrrrrrrr}
\toprule
$\mu$ & $\lambda$ & \multicolumn{3}{c}{Proposed} & \multicolumn{3}{c}{Brent} & \multicolumn{3}{c}{Floyd} \\
\cmidrule(lr){3-5}\cmidrule(lr){6-8}\cmidrule(lr){9-11}
 & & Next & Cmp & Total & Next & Cmp & Total & Next & Cmp & Total \\
\midrule
0 & 5 & 11 & 11 & 22 & 17 & 13 & 30 & 15 & 6 & 21 \\
10 & 7 & 35 & 47 & 82 & 49 & 33 & 82 & 62 & 25 & 87 \\
128 & 256 & 769 & 1023 & 1792 & 1023 & 640 & 1663 & 1024 & 385 & 1409 \\
256 & 1 & 515 & 1025 & 1540 & 1025 & 769 & 1794 & 1280 & 513 & 1793 \\
300 & 299 & 1199 & 1665 & 2864 & 1709 & 1111 & 2820 & 2394 & 899 & 3293 \\
\bottomrule
\end{tabular}
\end{table}

\begin{table}[H]\centering
\caption{Weighted cost $C_w=\#\text{next}+w\cdot\#\text{cmp}$: fraction of the
$90{,}300$ inputs on which the proposed algorithm has no larger weighted cost.}
\label{tab:weighted}
\begin{tabular}{rrrrr}
\toprule
$w$ & Proposed $\le$ Brent & Proposed $\le$ Floyd & Mean $A/B$ & Mean $A/F$ \\
\midrule
0.00 & 1.0000 & 1.0000 & 0.7230 & 0.6840 \\
0.05 & 1.0000 & 0.9285 & 0.7436 & 0.7125 \\
0.10 & 1.0000 & 0.8999 & 0.7629 & 0.7400 \\
0.25 & 1.0000 & 0.7645 & 0.8143 & 0.8168 \\
0.50 & 1.0000 & 0.4969 & 0.8829 & 0.9283 \\
1.00 & 0.4975 & 0.4529 & 0.9789 & 1.1051 \\
2.00 & 0.1229 & 0.1413 & 1.0888 & 1.3440 \\
\bottomrule
\end{tabular}
\end{table}

The weighted study (Table~\ref{tab:weighted}) makes the scope precise: the
proposed method is universally no worse than Brent when a comparison costs at most
half a successor evaluation ($w\le\tfrac12$), but the guarantee disappears as
comparisons approach the cost of successor evaluations ($w\to1$). This is exactly
the boundary that the wall-clock study below crosses.

\section{Empirical Evaluation: Wall-Clock and Expensive Oracles}
\label{sec:wallclock}

Operation counts alone do not establish wall-clock behaviour. We implemented all
three algorithms over a functional graph in which they interact with the data
\emph{only} through a successor primitive and integer-identity comparisons, with a
parameter \textsc{cost} controlling the price of one successor step:
$\textsc{cost}=0$ is a bare pointer load (the cheapest possible successor), while
larger \textsc{cost} injects per-step arithmetic modelling an expensive successor
such as a modular map. Measurements use $\mu=\lambda$ up to $2\times10^6$ on a
single core.

\begin{table}[H]\centering
\caption{Wall-clock time per run (ms) and successor-evaluation counts at
$\mu=\lambda=10^6$ ($n=2\times10^6$). Successor counts confirm the model:
the proposed method performs the fewest evaluations.}
\label{tab:wallclock}
\begin{tabular}{rrrrrrr}
\toprule
 & \multicolumn{3}{c}{Successor evaluations} & \multicolumn{3}{c}{Wall-clock (ms)} \\
\cmidrule(lr){2-4}\cmidrule(lr){5-7}
\textsc{cost} & Floyd & Brent & Proposed & Floyd & Brent & Proposed \\
\midrule
0  & 5{,}000{,}000 & 5{,}048{,}575 & 4{,}000{,}001 & 4.61 & 6.39 & 6.08 \\
8  & 5{,}000{,}000 & 5{,}048{,}575 & 4{,}000{,}001 & 27.59 & 29.14 & 24.04 \\
64 & 5{,}000{,}000 & 5{,}048{,}575 & 4{,}000{,}001 & 440.96 & 447.58 & 357.20 \\
\bottomrule
\end{tabular}
\end{table}

Two facts stand out. First, the model is faithful: the proposed algorithm really
does perform the fewest successor evaluations ($4.0\text{M}$ vs $5.0\text{M}$ for
Floyd and $5.05\text{M}$ for Brent), independent of \textsc{cost}. Second, fewer
evaluations do not imply less wall-clock time when a successor is cheap: at
$\textsc{cost}=0$, Floyd is fastest despite the most dereferences, because the
proposed method pays for its extra comparisons and richer control flow. Only once
a successor step is roughly an order of magnitude more expensive than a
comparison ($\textsc{cost}\gtrsim8$) does the proposed method become fastest ---
exactly the regime in which Brent was historically preferred over Floyd.

\paragraph{A tunable-cost scenario.}
We corroborate this with a deterministic iterated map
$f(x)=\mathrm{mix}^{K}(x)\bmod N$ ($N=2^{20}$, strong 64-bit mixing, salted per
instance), where $K$ tunes the cost of one oracle call. All finders use $O(1)$
memory; ground truth (a visited array) validates that every run recovers the true
$(\mu,\lambda)$. Over $1500$ random maps per cost level there were no validation
failures, the structural constant is cost-invariant ($3.46$ Floyd, $3.08$ Brent,
$2.18$ proposed), and the wall-clock speedup is present at every level, tending to
the evaluation ratio $3.08/2.18\approx1.41$ over Brent (and $\approx1.58$ over
Floyd) as the oracle grows expensive (Table~\ref{tab:scenario},
Figure~\ref{fig:scenario}).

\begin{table}[H]\centering
\caption{Period detection on a tunable-cost deterministic map ($1500$ maps/level).}
\label{tab:scenario}
\begin{tabular}{ccccc}
\toprule
$K$ & Floyd (ms) & Brent (ms) & Proposed (ms) & speedup vs Brent\\
\midrule
1  & 42.5 & 53.4 & 39.0 & $1.37\times$\\
4  & 86.8 & 123.0 & 79.2 & $1.55\times$\\
16 & 277.7 & 327.2 & 246.2 & $1.33\times$\\
32 & 668.8 & 667.3 & 487.1 & $1.37\times$\\
64 & 1431.0 & 1356.1 & 973.6 & $1.39\times$\\
\bottomrule
\end{tabular}
\end{table}

\begin{figure}[H]\centering
\includegraphics[width=0.92\linewidth]{figures_nextval/fig_scenario_prng.pdf}
\caption{(a) Cost-invariant evaluation constant. (b) The proposed method is
fastest at every oracle cost. (c) The wall-clock speedup approaches the
evaluation ratio as the oracle grows expensive.}
\label{fig:scenario}
\end{figure}

\section{Application: Pollard's $\rho$ Factorization}
\label{sec:rho}

The successor-evaluation model's canonical realization is Pollard's $\rho$
method~\cite{pollard1975}, where the successor is $f(x)=x^2+c \bmod N$, so that
\emph{one successor evaluation is exactly one modular squaring} --- the dominant
arithmetic cost --- while node equality is tested by the gcd signal
$1<\gcd(|x_i-x_j|,N)<N$. Factorization needs a \emph{collision}, not the entry,
so the relevant cost is detection: by Lemma~\ref{lem:brent-detect} the proposed
schedule never evaluates more successors than Brent to reach a collision.

\paragraph{Setup.}
We factored $400$ semiprimes $N=pq$ with $p\in[3\times10^4,3\times10^5]$ and
$q\in[3\times10^5,3\times10^6]$. For each $N$ we used the first additive constant
$c\in\{1,\dots,40\}$ for which all three strategies returned a non-trivial factor
on the same instance, and counted modular squarings until a factor was found.

\begin{table}[H]\centering
\caption{Pollard's $\rho$: modular squarings to find a factor, averaged over
$400$ semiprimes. The proposed schedule uses $17.7\%$ fewer squarings than Brent
and $47.0\%$ fewer than Floyd, and is never worse on any instance.}
\label{tab:rho}
\begin{tabular}{lccc}
\toprule
 & Proposed $A$ & Brent $B$ & Floyd $F$\\
\midrule
mean modular squarings & $541.3$ & $657.4$ & $1021.5$\\
mean ratio to proposed & $1.00$ & $0.859$ ($A/B$) & $0.544$ ($A/F$)\\
never-worse rate & --- & $400/400$ & $400/400$\\
strictly-better rate & --- & $166/400$ & $400/400$\\
\bottomrule
\end{tabular}
\end{table}

The measured advantage ($17.7\%$ over Brent, $47.0\%$ over Floyd) is consistent
with the operation-count predictions of Section~\ref{sec:empirical} and is
realized on a workload whose successor cost is intrinsically expensive. The ties
against Brent ($234/400$) occur exactly when the two detection schedules coincide
--- the phenomenon characterized in Theorem~\ref{thm:brent-ties-closed}. This is
the regime that originally motivated Brent over Floyd, and it is precisely where
the proposed method improves further.

\subsection{Hardened evaluation: 64-bit Montgomery arithmetic}
\label{sec:rho-hardened}

To rule out implementation artifacts we re-ran the comparison with a
production-style kernel: 64-bit modular arithmetic in Montgomery form
(\texttt{\_\_int128} REDC, no \texttt{libgmp} dependency), the map $f(x)=x^2+c$
counted as one Montgomery squaring per successor, batched gcd over windows of
$128$ accumulated differences, and deterministic Miller--Rabin prime generation.
For each instance the exact $(\mu_p,\lambda_p)$ modulo the smaller prime $p$ is
obtained by an independent Brent walk, and squaring counts are then taken from the
validated closed forms, so the comparison is exact rather than sampled.

\begin{table}[H]\centering
\caption{Modular squarings to the first non-trivial factor, $400$ random balanced
semiprimes $N=pq$ per bit-size. The proposed schedule needs $\approx30\%$ fewer
squarings than Brent and $\approx37\%$ fewer than Floyd, and is never worse on any
single instance.}
\label{tab:rho-hardened}
\begin{tabular}{rrrrccc}
\toprule
$\log_2 N$ & $\mathbb E[A]$ & $\mathbb E[B]$ & $\mathbb E[F]$ & $A/B$ & $A/F$ & strict $<$ both\\
\midrule
32 & 629.1 & 894.9 & 988.8 & 0.703 & 0.636 & $100\%$\\
40 & 2{,}362 & 3{,}380 & 3{,}917 & 0.699 & 0.603 & $100\%$\\
48 & 9{,}808 & 13{,}895 & 15{,}475 & 0.706 & 0.634 & $100\%$\\
56 & 40{,}875 & 57{,}799 & 63{,}916 & 0.707 & 0.640 & $100\%$\\
60 & 77{,}805 & 110{,}839 & 127{,}157 & 0.702 & 0.612 & $100\%$\\
\bottomrule
\end{tabular}
\end{table}

The ratios are flat across four orders of magnitude in $N$, exactly as the
constant-factor theory predicts. Wall-clock timing on the same kernel
(Table~\ref{tab:rho-wallclock}) confirms the saving survives end-to-end: adding a
per-step ``expensive-oracle'' load \textsc{extra} (modelling co-factor work or a
heavier mixing map), the proposed method is $\approx1.43\times$ faster than Brent
and $\approx1.6\times$ faster than Floyd at every level.

\begin{table}[H]\centering
\caption{Wall-clock per factored $48$-bit semiprime ($\mu$s, mean$\pm$s.d.,
$400$ instances/level; Montgomery REDC, batched gcd). \textsc{extra} is the
injected per-step oracle cost.}
\label{tab:rho-wallclock}
\begin{tabular}{rrrrcc}
\toprule
\textsc{extra} & Floyd & Brent & Proposed & vs Brent & vs Floyd\\
\midrule
0  & $150.2\pm84.2$ & $132.1\pm66.7$ & $92.3\pm45.4$ & $1.43\times$ & $1.63\times$\\
2  & $257.4\pm144$ & $226.6\pm113$ & $158.2\pm77$ & $1.43\times$ & $1.63\times$\\
8  & $629.8\pm500$ & $528.8\pm266$ & $369.7\pm181$ & $1.43\times$ & $1.70\times$\\
32 & $2064\pm1489$ & $1761\pm901$ & $1231\pm647$ & $1.43\times$ & $1.68\times$\\
64 & $3805\pm2132$ & $3360\pm1696$ & $2376\pm1213$ & $1.41\times$ & $1.60\times$\\
\bottomrule
\end{tabular}
\end{table}

\section{Related Work}
\label{sec:related}

Floyd's tortoise-and-hare is commonly attributed to his 1967 work on
nondeterministic algorithms~\cite{floyd1967}; Brent's power-of-two method arose
in Monte-Carlo factorization with a reported one-third reduction over Floyd in
Pollard-$\rho$ applications~\cite{brent1980}, which our measurements reproduce
(Section~\ref{sec:rho}). Pollard's $\rho$~\cite{pollard1975} is the canonical
setting where a successor is genuinely expensive. The distribution of
$(\mu,\lambda)$ for random mappings is governed by the Knuth--Flajolet
school~\cite{flajolet,sedgewick2013}, with $\mathbb E[\mu+\lambda]\sim\sqrt{\pi/2}\sqrt n$.
A second line trades space for evaluations: Gosper's stack scheme~\cite{gosper1972}
and Nivasch's single-pass stack~\cite{nivasch2004} reach the optimum
$\mu+\lambda$ but use $\Theta(\log(\mu+\lambda))$ memory; the time--space bounds
of Sedgewick, Szymanski and Yao~\cite{ssy1982} show this memory is necessary.
Table~\ref{tab:related} places the present method, using the closed forms
established above; here $k_B=\lceil\log_2\max(\lambda,\mu+1)\rceil$ and
$T=\lambda\lceil\mu/\lambda\rceil$.

\begin{table}[H]\centering
\caption{Cycle-detection methods by auxiliary space and successor-evaluation
count. The proposed method dominates both constant-space baselines pointwise
while remaining in $O(1)$ space; Nivasch attains the evaluation optimum but uses
non-constant memory.}
\label{tab:related}
\begin{tabular}{llll}
\toprule
Method & Aux.\ space & Successor evaluations (cyclic) & Notes\\
\midrule
Floyd~\cite{floyd1967} & $O(1)$ & $F=3T+2\mu$ & tortoise/hare\\
Brent~\cite{brent1980} & $O(1)$ & $B=2^{k_B}\!-\!1+2\lambda+2\mu$ & power-of-two\\
\textbf{Proposed} & $O(1)$ & $\mu+\lambda\le A\le 2\mu+3\lambda$, $A\le\min(B,F)$ & this paper\\
Nivasch~\cite{nivasch2004} & $O(\log k)$ exp. & $\mu+\lambda$ (optimal) & stack\\
\midrule
Lower bound (Thm.~\ref{thm:lowerbound}) & any & $\ge\mu+\lambda$ & this paper\\
\bottomrule
\end{tabular}
\end{table}

\paragraph{Expensive successors and parallel search.}
Where the successor is a cryptographic map, three further lines are relevant. The
Pollard--Brent factorization of $F_8$~\cite{brent1981} established the $\rho$
method at scale, and Montgomery's batching and inversion
tricks~\cite{montgomery1987} reduce the per-step cost we hold fixed: they are
orthogonal to, and compose with, our schedule, which changes only \emph{how many}
successors are taken. Teske's adding- and mixing-walks~\cite{teske1998} improve
the statistical quality of the iterating map (shrinking the effective
$\mu+\lambda$), again orthogonally to the detection schedule. Finally, the
distinguished-point technique of Quisquater and Delescaille~\cite{qd1990} and the
parallel collision search of van Oorschot and Wiener~\cite{vow1999} trade memory
and many machines for throughput; our contribution is to the \emph{serial}
constant of a single walk, and a distinguished-point variant of the localization
anchor is a natural extension (Section~\ref{sec:limits}).

Our contribution is not a new asymptotic class but a tighter constant with an
exact, verifiable analysis, a proved per-instance dominance, a matching lower
bound, and a multi-anchor latency/work tradeoff.

\section{Discussion and Limitations}
\label{sec:limits}
\begin{itemize}[leftmargin=2em,itemsep=2pt]
  \item The asymptotic complexity remains $\Theta(\mu+\lambda)$; the improvement
        is a constant factor ($\approx2.18$ in expectation versus $3.08$/$3.5$).
  \item The dominance results count only successor evaluations on cyclic inputs.
        The method performs more pointer comparisons than Floyd and has more
        control flow than Brent; in the cheap-successor regime these costs make
        it slower than plain Floyd (Section~\ref{sec:wallclock},
        Table~\ref{tab:weighted}).
  \item For acyclic lists, whether a null check counts as a dereference changes
        exact constants, so we claim dereference-count dominance only for cyclic
        inputs and only $O(n)$ termination otherwise.
  \item In some applications (e.g.\ Pollard $\rho$) other costs --- batched gcds,
        modular reductions --- may matter, so the end-to-end saving must be
        re-measured per application.
  \item The method \emph{detects} recurrence; it does not by itself \emph{escape}
        it. In optimization-style settings it is a cheap oscillation/limit-cycle
        trigger, not a solver.
  \item \textbf{Cost model.} The dominance metric counts successor evaluations
        only. On hardware where comparisons and branch mispredictions are not
        free, a net speedup requires the successor to cost at least
        $\approx8\times$ a comparison (Table~\ref{tab:weighted},
        Section~\ref{sec:wallclock}); we state this threshold explicitly rather
        than claim an unconditional speedup.
  \item \textbf{Originality.} The two-phase probe and power-of-two anchors are
        Brent-style; what is new, to our knowledge, is combining a
        \emph{failure-compression} (nextval) prefix-skip with a
        \emph{localization anchor} (the previous round's anchor, proved safe to
        start entry detection from). The skip is what lowers \emph{total} --- not
        merely detection --- evaluations, yielding the $\approx2.18$ constant that
        pure multi-anchor schemes provably cannot reach
        (Corollary~\ref{cor:o1-optimal}).
  \item \textbf{Open dimensions.} A distinguished-point/parallel variant,
        SIMD-vectorised batching of the squaring, online/streaming inputs, and
        robustness to a noisy successor oracle are untouched here, as is closing
        the gap between the proved $O(1)$-space detection floor $1.2213$ and the
        total constant $2.18$.
\end{itemize}

\paragraph{Positioning.}
The honest summary is that this is a constant-factor refinement, not a
complexity-class result: $\Theta(\mu+\lambda)$ time and $O(1)$ space are
unchanged. Its value concentrates where (i) the successor dominates cost
(Pollard $\rho$, expensive mixing maps) and (ii) per-instance worst-case
guarantees matter, since the method never loses to Brent or Floyd on any cyclic
input and wins on almost all.

\section{Conclusion}
\label{sec:conclusion}

We presented a monotone single-pointer anchor-advancement algorithm for cycle
detection that combines two-phase probing, a failure-compression (nextval) skip
rule, and a localization anchor that starts entry detection closer to the cycle.
We proved correctness, linear time, and constant space; gave an exact closed form
for its successor-evaluation count; proved that it weakly dominates both Brent and
Floyd on \emph{every} cyclic input --- strictly so on all but three (resp.\ one)
degenerate inputs --- and that, against a matching $\mu+\lambda$ lower bound, it is
constant-factor optimal. An average-case Mellin analysis identifies its constant
as a log-periodic $\approx2.18033$, a multi-anchor extension traces the
detection/total-work Pareto frontier, and wall-clock and Pollard-$\rho$ studies
locate the crossover at which fewer evaluations become a real speedup. The result
is a clean, distribution-free, fully analyzable refinement of Brent-style cycle
finding in the successor-evaluation model --- most useful wherever the successor
is the dominant cost --- rather than a universal replacement.

\begin{thebibliography}{9}

\bibitem{floyd1967} R.~W.~Floyd. Nondeterministic algorithms.
\emph{Journal of the ACM}, 14(4):636--644, 1967.

\bibitem{brent1980} R.~P.~Brent. An improved Monte Carlo factorization algorithm.
\emph{BIT Numerical Mathematics}, 20:176--184, 1980.

\bibitem{pollard1975} J.~M.~Pollard. A Monte Carlo method for factorization.
\emph{BIT Numerical Mathematics}, 15(3):331--334, 1975.

\bibitem{kmp1977} D.~E.~Knuth, J.~H.~Morris~Jr., and V.~R.~Pratt. Fast pattern
matching in strings. \emph{SIAM Journal on Computing}, 6(2):323--350, 1977.

\bibitem{flajolet} P.~Flajolet and A.~M.~Odlyzko. Random mapping statistics.
\emph{EUROCRYPT '89}, LNCS 434, 329--354, 1990.

\bibitem{nivasch2004} G.~Nivasch. Cycle detection using a stack.
\emph{Information Processing Letters}, 90(3):135--140, 2004.

\bibitem{gosper1972} R.~W.~Gosper. Item 132: cycle detection in $O(\log)$ space.
In M.~Beeler, R.~W.~Gosper, and R.~Schroeppel, \emph{HAKMEM}, MIT AI Memo 239, 1972.

\bibitem{ssy1982} R.~Sedgewick, T.~G.~Szymanski, and A.~C.~Yao. The complexity of
finding cycles in periodic functions. \emph{SIAM Journal on Computing},
11(2):376--390, 1982.

\bibitem{sedgewick2013} R.~Sedgewick and P.~Flajolet. \emph{An Introduction to
the Analysis of Algorithms}. Addison-Wesley, 2nd edition, 2013.

\bibitem{montgomery1987} P.~L.~Montgomery. Speeding the Pollard and elliptic curve
methods of factorization. \emph{Mathematics of Computation}, 48(177):243--264, 1987.

\bibitem{brent1981} R.~P.~Brent and J.~M.~Pollard. Factorization of the eighth
Fermat number. \emph{Mathematics of Computation}, 36(154):627--630, 1981.

\bibitem{teske1998} E.~Teske. Speeding up Pollard's rho method for computing
discrete logarithms. \emph{ANTS-III}, LNCS 1423, 541--554, 1998.

\bibitem{qd1990} J.-J.~Quisquater and J.-P.~Delescaille. How easy is collision
search? Application to DES. \emph{EUROCRYPT '89}, LNCS 434, 429--434, 1990.

\bibitem{vow1999} P.~C.~van Oorschot and M.~J.~Wiener. Parallel collision search
with cryptanalytic applications. \emph{Journal of Cryptology}, 12(1):1--28, 1999.

\end{thebibliography}

\end{document}
